% ============================================================
%  PLANTILLA DE MAPA MENTAL — syntheca
%  Clase STANDALONE/ARTICLE, usa la librería `mindmap` de TikZ
%  (distribución radial automática, no hay que calcular ángulos
%  a mano por nodo).
%
%  Filosofía de color, distinta al resto del sistema: acá SÍ se
%  usa color para diferenciar ramas (familias de conceptos) —
%  no es decoración, es una ayuda de codificación dual (dual
%  coding) para la memoria, la misma lógica por la que
%  pedagogia-cognitiva permite diagramas cuando cumplen una
%  función real. La paleta es apagada/profesional, nunca
%  saturada — mantiene el mismo espíritu sobrio del resto.
%
%  REGENERACIÓN COMPLETA: a diferencia del libro (que acumula
%  capítulos), este archivo se reescribe ENTERO cada vez que
%  generador-mapa-mental corre — nunca se edita a mano un nodo
%  suelto. La fuente de verdad es mapa-estudio.json, no el .tex.
%
%  ------------------------------------------------------------
%  COPIA LOCAL DEL REPO. El original vive en el plugin, en
%  ~/.claude/skills/syntheca/skills/plantilla-sintesis/latex/.
%  Esta copia es la que se edita; el plugin NO se toca.
%
%  CAMBIOS DE ESTA COPIA respecto de la del plugin (2026-09-05,
%  al pasar el mapa del Parcial I de 33 a 54 conceptos):
%    1. Tercer nivel de jerarquía (`level 3 concept`), para que
%       el mapa sea raíz -> familia -> subtema -> concepto y no
%       raíz -> familia -> concepto.
%    2. Distancias y anchos de texto mayores en todos los
%       niveles (era el fix ya identificado y nunca aplicado).
%    3. Lienzo configurable con \setLienzo: el diámetro natural
%       de un mapa radial crece de forma lineal con la cantidad
%       de hojas, así que el papel no puede ser fijo.
%    4. Se eliminó `\tikzstyle{every node}=[font=\sffamily]`.
%       Era un BUG: `every node` se aplica DESPUÉS de los
%       estilos de nivel, así que pisaba el `font=` de
%       `level 1/2/3 concept` y todos los nodos salían en
%       \normalsize. La familia sans ya la garantiza
%       \familydefault, así que borrarlo no cambia la tipografía
%       pero devuelve el control del CUERPO de letra por nivel,
%       que es de donde sale casi toda la capacidad de empaque.
%  ============================================================
\usepackage[utf8]{inputenc}
\usepackage{lmodern}   % antes de fontenc, para no caer en fuentes bitmap (Type3)
\usepackage[T1]{fontenc}
\usepackage{tikz}
\usetikzlibrary{mindmap,trees,backgrounds}
\usepackage{adjustbox}   % para encajar el mapa en la página (ver "Escalado")
\usepackage{fancyhdr}
\usepackage[table]{xcolor}
\usepackage{geometry}
\geometry{top=1.2cm, bottom=1.2cm, left=1cm, right=1cm, landscape}

% ---------- Lienzo ----------
% Un mapa radial tiene radio proporcional a la cantidad de hojas (todas las
% hojas viven sobre una misma circunferencia, así que el perímetro tiene que
% alcanzar para todas). Su ÁREA, entonces, crece con el cuadrado de la
% cantidad de conceptos: pasar de 17 hojas a 54 no triplica el mapa, lo
% multiplica por diez. Por eso el papel no puede ser fijo: con A4 apaisado
% un mapa de 54 conceptos sólo entra encogido a ~0.27, y a esa escala la
% letra queda en 2 pt. \setLienzo permite que cada mapa declare el papel que
% su cantidad de conceptos necesita, manteniendo el cuerpo de letra legible.
% Se llama en el cuerpo, antes de \begin{document}? No: geometry admite
% reconfiguración, así que se llama justo después de cargar el preámbulo.
\newcommand{\setLienzo}[2]{\geometry{paperwidth=#1, paperheight=#2, landscape=false,
  top=1.2cm, bottom=1.2cm, left=1.2cm, right=1.2cm}}

% Misma familia sans que el libro.
\renewcommand{\familydefault}{\sfdefault}

\definecolor{textgray}{RGB}{60, 60, 60}
\definecolor{rulegray}{RGB}{150, 150, 150}

% ---------- Paleta de ramas — apagada, profesional, máximo 6 colores.
% Si un examen tiene más de 6 familias, se reciclan en orden. ----------
\definecolor{ramaUno}{RGB}{91, 110, 130}    % azul-gris apagado
\definecolor{ramaDos}{RGB}{120, 100, 90}    % marrón-gris apagado
\definecolor{ramaTres}{RGB}{95, 120, 100}   % verde-gris apagado
\definecolor{ramaCuatro}{RGB}{115, 95, 120} % violeta-gris apagado
\definecolor{ramaCinco}{RGB}{130, 110, 80}  % ocre apagado
\definecolor{ramaSeis}{RGB}{90, 115, 120}   % teal-gris apagado

\newcommand{\miNombre}{Nombre Apellido}
\newcommand{\materiaActual}{Materia}
\newcommand{\setMateria}[1]{\renewcommand{\materiaActual}{#1}}
\newcommand{\nombreMapa}{Mapa mental}
\newcommand{\setNombreMapa}[1]{\renewcommand{\nombreMapa}{#1}}

\pagestyle{fancy}
\fancyhf{}
\renewcommand{\headrulewidth}{0pt}
\renewcommand{\footrulewidth}{0pt}
\fancyfoot[C]{\small\color{textgray}\thepage}
\fancyfoot[R]{\footnotesize\color{textgray}\miNombre}
\fancyfoot[L]{\footnotesize\color{textgray}\materiaActual}

% ---------- Estilo base de los conceptos (nodos) ----------
% Nivel raíz: el examen/tema central.
% Nivel 1: familias de conceptos (ramas, coloreadas).
% Nivel 2: subtemas dentro de una familia (heredan el color de la rama).
% Nivel 3: conceptos individuales (hojas, color heredado de su rama).
%
% El nivel 2 es NUEVO. Antes las hojas colgaban directo de la familia y una
% familia de 17 conceptos abría un abanico de 800°, imposible de acomodar.
% Con un nivel intermedio la misma familia se reparte en 2 o 3 abanicos
% chicos, y el ancho angular que necesita cae a menos de la mitad.
%
% grow cyclic es obligatorio en los tres niveles: sin él las ramas no se
% distribuyen en círculo, crecen todas hacia abajo y se apilan.
%
% text width fija el ancho de párrafo dentro del círculo; sin él, una
% etiqueta larga se desborda del nodo o lo infla de forma desproporcionada.
% minimum size va un poco por encima de text width para que quede aire.
%
% Las distancias de abajo NO son decorativas: salieron de una búsqueda sobre
% la grilla (level distance x sibling angle) verificando, para cada
% combinación, que no hubiera ni un par de nodos con distancia entre centros
% menor que la suma de sus radios. Cambiar una a ojo rompe esa garantía.
\tikzset{
  root concept/.style={
    concept color=textgray,
    font=\sffamily\bfseries\Large,
    text=white,
    minimum size=4.4cm,
    text width=3.6cm,
  },
  level 1 concept/.append style={
    grow cyclic,
    font=\sffamily\bfseries\small,
    text=white,
    level distance=19cm,
    minimum size=3.0cm,
    text width=2.5cm,
  },
  level 2 concept/.append style={
    grow cyclic,
    font=\sffamily\bfseries\footnotesize,
    text=white,
    sibling angle=52,
    level distance=7.5cm,
    minimum size=2.45cm,
    text width=2.0cm,
  },
  level 3 concept/.append style={
    grow cyclic,
    font=\sffamily\scriptsize,
    text=white,
    % La cuerda entre dos hojas hermanas es 2*d*sin(a/2); con d=4.6cm y a=34
    % da 2.69cm, contra un diámetro de hoja de ~2.13cm: sobran 0.56cm. Ese
    % número es la holgura mínima de TODO el mapa (ningún otro par queda más
    % cerca), así que bajar `sibling angle` o `level distance` acá es lo
    % primero que lo rompe.
    sibling angle=34,
    level distance=4.6cm,
    minimum size=2.0cm,
    text width=1.7cm,
  },
}

% ---------- Número de sectores angulares ----------
% El ángulo entre hijos de la raíz es 360/n. Con \setRamas{<familias>} todas
% las ramas reciben la MISMA porción, lo que sólo sirve si todas tienen una
% cantidad parecida de hojas. Cuando no la tienen (Parcial I: familias de 2 y
% de 17 conceptos), conviene repartir el círculo en SLOTS iguales y darle a
% cada rama varios slots, proporcionalmente a su cantidad de hojas: se llama
% \setRamas{<slots totales>} y en el cuerpo se rellena con `child[missing]`,
% que consume un slot sin dibujar nada. \setSlots es el mismo comando con el
% nombre que corresponde a ese uso.
\newcommand{\setRamas}[1]{%
  \tikzset{level 1 concept/.append style={sibling angle=360/#1}}%
}
\let\setSlots\setRamas
\setRamas{4}

% ---------- Escalado ----------
% Para que el mapa entre en la página, envolver TODO el tikzpicture en el
% entorno \begin{mapaAjustado} ... \end{mapaAjustado}, que lo encoge sólo si
% hace falta para respetar ancho Y alto disponibles. NUNCA usar scale=...
% junto con transform shape: grow cyclic rota el sistema de coordenadas de
% cada rama y transform shape le aplica esa rotación también al texto, así
% que la mitad de las etiquetas salen cabeza abajo.
%
% Ojo: `mapaAjustado` es una red de seguridad, no un plan. Si el mapa entra
% sólo gracias a un encogido fuerte, el problema es el lienzo (\setLienzo),
% porque el encogido se lleva puesto el cuerpo de letra.
\newenvironment{mapaAjustado}{%
  \begin{center}%
  \begin{adjustbox}{max size={\linewidth}{0.88\textheight}}%
}{%
  \end{adjustbox}%
  \end{center}%
}
